\subsection{\texorpdfstring{\(\mathbf{R}^n\) 的拓扑}{R 的 n 次幂的拓扑}}

定义 1. 与实直线相同的 n 个空间的拓扑积称为 n 维实数空间，记作 \(R^{n}\)。

注. 空间 \(R^{0}\) 由单独一个点组成。

由《集合论》第 III 章，§ 6，第 3 小节，定理 2 的推论 1 可知：若 E 是无限集，则 \(E^{n}\) 对一切整数 n \textgreater{} 0 与 E 等势；因此，若 n \textgreater{} 0 ，则 \(R^{n}\) 与 R 等势，也就是说，\(R^{n}\) 具有连续统的势（参见习题 1 与习题 2）。

定义 2. \(\mathbf{R}^n\) 中由 \(n\) 个 \(\mathbf{R}\) 的开（相应地，闭）区间之积构成的任何子集，称为 \(\mathbf{R}^n\) 中的开（相应地，闭）长方体。{[}当 \(n = 2\) 时，它称为开（相应地，闭）矩形。{]}

\(\mathbf{R}^n\) 中的开长方体构成 \(\mathbf{R}^n\) 的拓扑的一个基（第 I 章，§ 4，第 1 小节）；包含点 \(x = (x_i)_{1 \leqslant i \leqslant n}\) （\(\mathbf{R}^n\) 中的）的开长方体构成 \(x\) 的一个基本邻域系，而 \(\mathbf{R}^n\) 的以 \(x\) 为内点的闭长方体亦然。

\(\mathbf{R}^n\) 中每个非空开长方体同胚于 \(\mathbf{R}^n\)（第 IV 章，§ 4，第 1 小节，命题 1）。

由此推出：当 \(n \geqslant 1\) 时，\(R^{n}\) 中每个非空开集具有连续统的势。

\(R^{n}\) 的开（相应地，闭）立方体是指 n 个等长有界区间之积所成的开（相应地，闭）长方体{[}当 n = 2 时，它称为开（相应地，闭）正方形{]}；这些区间的公共长度称为立方体的边（或边长）。开立方体

\[
\mathrm{K} _ {m} = \prod_ {1 \leqslant i \leqslant n} \left] x _ {i} - \frac {1}{m}, x _ {i} + \frac {1}{m} \right[
\]

构成点 \(\boldsymbol{x} = (x_{i})\) 的可数基本邻域系，这里 m 取遍全体整数 \textgreater0 的集，或取遍任一递增趋于无穷的整数序列。

\(\mathbf{R}^n\) 中每个开（或闭）长方体都是连通的（第 I 章，§ 11，第 4 小节，命题 8）；特别地，\(\mathbf{R}^n\) 是连通的且局部连通的。

若 A 是 \(R^{n}\) 中的非空开集，则它的分支都是开集（第 I 章，§ 11，第 6 小节，命题 11）；而且这些分支组成的集是可数的，因为 \(R^{n}\) 有可数稠密集（例如 \(Q^{n}\)）。

考察 \(R^{n}\) 的子集 A 在什么条件下是相对紧的。由吉洪诺夫定理（第 I 章，§ 9，第 5 小节，定理 3），必要且充分的是 A 在 \(R^{n}\) 的诸因子上的投影都是相对紧的；由波莱尔–勒贝格定理（第 IV 章，§ 2，第 2 小节，定理 2），这等价于说这些投影都是 R 的有界子集。我们说 \(R^{n}\) 的子集 A 是有界的，如果它的所有投影都是 R 的有界子集；于是我们证明了：

命题 1. \(R^{n}\) 的子集 A 相对紧，当且仅当它是有界的。

推论. 空间 \(\mathbf{R}^n\) 是局部紧的，但当 \(n \geqslant 1\) 时不是紧的。

